\itemhead{第1章 小结与复习 · 第1题（教材p.28）}{XJ-B1-C01-SU-P028-Y0974}
\label{XJ-B1-C01-SU-P028-Y0974}
{\footnotesize 湘教版教材p.28 / PDF第35页 / 第1章 小结与复习 / 第1题}
\field{原题}{集合是现代数学的基本概念。用集合的语言和符号表达数学事实，比起自然语言更为严谨和简洁，而且世界通用。试结合生活中、数学中以及其他学科中的集合实例，用集合语言来描述它们，并交流用集合语言来表达的优势。}
\field{参考答案}{示例：全班学生组成集合 $U$，戴眼镜者组成 $A$，则“某同学戴眼镜”可写为 $x\in A$。集合语言能明确对象范围、关系和运算，避免自然语言歧义。}
{\small \field{关联知识}{S01 集合、元素与常用数集（教材帮p.1、2 / PDF10、11）；S02 集合表示、区间与数集点集（教材帮p.2、3、4 / PDF11、12、13）}
\field{方法与易错点}{先辨认研究对象和数集范围，再代入定义判断；含参集合逐个来源分类。 易错：把自然数集误作正整数集；混淆元素和单元素集合；候选参数造成元素重复。}
\field{答案依据}{答案册未提供本栏目的独立对应答案；依据原题、教材概念补写参考回应。开放性总结不设唯一数值答案。}}
